\begin{korollar}{Gordon'sche Paare und normiterierende Mengen}
{GordonschePaareUndNormiterierendeMengen}
Seien $(E, \norm\cdot_E)$ ein normierter $\R$--Vektorraum, $x\in E^N$ und
$G_x$ der Gauß'sche Operator zum Tupel $(x,g)$.\\
Seien $\emptyset \neq T \subseteq S^{m-1}$,
$\emptyset\neq V\subseteq \R^N$ und $c \in \R_{> 0}$ mit
$\forall\, v \in V: \norm v_{2, N} \le c$.\\
Es gelte:
\begin{enumerate}[(a)]
	\item $0 \in V \folgt 0$ Häufungspunkt von $V$.
	\item $\forall\, w \in \R^N: \sup\menge{\sprod w v _N}{v \in V} 
	= \norm{\bprod w x_N}_E$.
\end{enumerate}
Dann sind nach \refclaim{IntegrierbarkeitVonExtremaInGordonschenPaaren}
{IntegrierbarkeitVonExtremaInGordonschenPaaren1}
\[\norm\cdot_E \circ \bprod \cdot x _N \circ \gamma, \quad
\inf_{t\in T}\bprod t \delta_m, \quad \sup_{t\in T}\bprod t \delta_m\]
und nach \refclaim{IntegrierbarkeitVonExtremaInGordonschenPaaren}
{IntegrierbarkeitVonExtremaInGordonschenPaaren2}
\[\inf_{t \in T}\, \norm\cdot_E \circ ev(t,\cdot)\circ G_x, \quad
\sup_{t \in T}\, \norm\cdot_E \circ ev(t,\cdot)\circ G_x \]
integrierbar und es gilt:
\begin{enumerate}[(1)]
	\item \label{GordonschePaareUndNormiterierendeMengen1}
	\begin{align*}
	& \E\left( \norm\cdot_E \circ \bprod \cdot x _N \circ \gamma\right)
	+ c \E\left(\inf_{t\in T}\bprod t \delta_m \right)\\
	\le &\,\E\left(\inf_{t \in T} \norm\cdot_E \circ ev(t,\cdot)\circ G_x \right)
	\le \E\left(\sup_{t \in T} \norm\cdot_E \circ ev(t,\cdot)\circ G_x \right)\\
	\le &\,\E\left( \norm\cdot_E \circ \bprod \cdot x _N \circ \gamma\right)
	+ c \E\left(\sup_{t\in T}\bprod t \delta_m \right)\text.
	\end{align*}
	\item \label{GordonschePaareUndNormiterierendeMengen2}
	Ist $T = -T$, so ist
	$\E\left(\Inf{t\in T}\bprod t \delta_m \right)
	= - \E\left(\Sup{t\in T}\bprod t \delta_m \right)$.
\end{enumerate}
\end{korollar}
\begin{proof}
(1) Sei $(X,Y)$ das Gordon'sche Paar zum Tupel $(\gamma,\delta,g,c)$.\\
Dann sind die Voraussetzungen von
\ref{GordonschePaareUndDerGaussscheMinMaxSatz} erfüllt.\\
Ferner gilt für alle $\omega \in \Omega$:
\begin{align}
\left(\sup_{v \in V}\bprod v \gamma_N \right)(\omega)
&= \sup_{v \in V}\sprod v {\gamma(\omega)_N}
=\sup_{v \in V}\sprod{\gamma(\omega)}v_N
\overset{\text{Vor.}}= \norm{\bprod {\gamma(\omega)} x}_N\notag\\
&= \left(\norm\cdot_E \circ \bprod \cdot x_N \circ \gamma\right) (\omega)\text,
\label{gpunmeq:1}
\end{align}
sowie für alle $s \in T$:
\begin{align}
\left(\sup_{v \in V}Y(s,v)\right)(\omega)&= \sup_{v \in V}\sprod{g(\omega)s}v_N
\overset{\text{Vor.}}= \norm{\bprod {g(\omega)s} x_N}_E
= \norm{G_x(\omega)s}_E \notag\\
&=\left(\norm\cdot_E \circ ev(s,\cdot)\circ G_x\right)(\omega)\text.
\label{gpunmeq:2}
\end{align}
Somit gilt wegen der Integrierbarkeit aller vorkommenden Abbildungen:
\begin{align*}
&\E\left( \norm\cdot_E \circ \bprod \cdot x _N \circ \gamma\right)
	+ c \E\left(\inf_{t\in T}\bprod t \delta_m \right)\\
	\overset{c>0}=&\,
	\E\left( \norm\cdot_E \circ \bprod \cdot x _N \circ \gamma\right)
	+  \E\left(\inf_{t\in T}c\bprod t \delta_m \right)\\
\overset{\eqref{gpunmeq:1}}= &\, \E\left(\sup_{v \in V}\bprod v \gamma_N
+ \inf_{t\in T}c\bprod t \delta_m\right)\\
\overset{
\refclaim{BeschraenktheitGordonscherPaare}{BeschraenktheitGordonscherPaare3}}=&
\,\E \left(\inf_{t \in T}\sup_{v \in V}X(t, v) \right)
\overset{\refclaim{GordonschePaareUndDerGaussscheMinMaxSatz}
{GordonschePaareUndDerGaussscheMinMaxSatz1}}
\le \E\left(\inf_{t \in T}\sup_{v \in V}Y(t, v) \right)\\
\overset{\eqref{gpunmeq:2}}= & \,
\E\left(\inf_{t \in T}\norm\cdot_E \circ ev(t,\cdot) \circ G_x \right)
\le
\E\left(\sup_{t \in T}\norm\cdot_E \circ ev(t,\cdot) \circ G_x \right)\\
\overset{\eqref{gpunmeq:1}}= 
&\, \E\left(\sup_{t \in T}\sup_{v \in V}Y(t, v) \right)
\overset{\refclaim{GordonschePaareUndDerGaussscheMinMaxSatz}
{GordonschePaareUndDerGaussscheMinMaxSatz2}}\le
\E\left(\sup_{t \in T}\sup_{v \in V}X(t, v) \right)\\
\overset{\refclaim{BeschraenktheitGordonscherPaare}
{BeschraenktheitGordonscherPaare3}} = &\,
\E\left(\sup_{v \in V}\bprod v \gamma_N
+ \sup_{t\in T}c\bprod t \delta_m\right)\\
\overset{\eqref{gpunmeq:1}}=&\,
\E\left( \norm\cdot_E \circ \bprod \cdot x _N \circ \gamma\right)
	+  \E\left(\sup_{t\in T}c\bprod t \delta_m \right)
\end{align*}
(2) Es gelte $T=-T$ und sei $\omega \in \Omega$. Dann gilt:
\begin{align*}
\left(\inf_{t\in T}\bprod t \delta_m\right)(\omega)
&= \inf_{t \in T}\sprod t {\delta(\omega)}_m
= - \sup_{t\in T}\left(-\sprod t {\delta(\omega)}\right)
= - \sup_{t\in T}\left(\sprod {-t} {\delta(\omega)}\right)\\
&=- \sup_{t\in -T}\left(\sprod t {\delta(\omega)}\right)
=- \sup_{t\in T}\left(\sprod t {\delta(\omega)}\right)
= \left(- \sup_{t\in T}\bprod t \delta_m \right)(\omega)\text.
\end{align*}
Da die genannten Abbildungen integrierbar sind, gilt somit (2).
\end{proof}